\documentclass[12pt, a4paper, openright]{report}
% openright: ogni capitolo inizia su pagina dispari (destra), come da linee guida

\usepackage[utf8]{inputenc}
\usepackage[T1]{fontenc}
\usepackage[italian]{babel}
\usepackage{geometry}
\usepackage{hyperref}
\usepackage{titlesec}
\usepackage{mathptmx}   % Times New Roman (consigliato dalle linee guida)
\usepackage{setspace}   % Gestione interlinea
\usepackage{microtype}  % Migliore giustificazione del testo

% -------------------------------------------------------
% MARGINI: tutti 3 cm (linee guida: 3 sx, 3 dx, 3 sup, 3 inf)
% -------------------------------------------------------
\geometry{top=3cm, bottom=3cm, left=3cm, right=3cm}

% -------------------------------------------------------
% INTERLINEA 1,5 per il corpo del testo
% NB: il pacchetto setspace mantiene automaticamente
%     le note a piè di pagina a interlinea singola
% -------------------------------------------------------
\onehalfspacing

% -------------------------------------------------------
% NOTE A PIÈ DI PAGINA: 10pt, giustificate, interlinea singola
% (setspace gestisce già la singola spaziatura automaticamente)
% -------------------------------------------------------
\renewcommand{\footnotesize}{\fontsize{10pt}{12pt}\selectfont}

% -------------------------------------------------------
% FORMATO TITOLI CAPITOLI
% "Capitolo N" su una riga, titolo sulla riga successiva
% -------------------------------------------------------
\titleformat{\chapter}[display]
  {\normalfont\bfseries\large}
  {\chaptername\ \thechapter}
  {8pt}
  {\Large}

\titleformat{\section}
  {\normalfont\bfseries\normalsize}
  {\thesection.}
  {0.5em}
  {}

\titleformat{\subsection}
  {\normalfont\itshape\normalsize}
  {\thesubsection.}
  {0.5em}
  {}

% -------------------------------------------------------
% HYPERREF: link in nero per la stampa
% -------------------------------------------------------
\hypersetup{
    colorlinks=true,
    linkcolor=black,
    filecolor=black,
    urlcolor=black,
    pdftitle={Titolo della Tesi},
    pdfauthor={Nome Cognome},
}


% =======================================================
\begin{document}
% =======================================================

% -------------------------------------------------------
% FRONTESPIZIO (template da linee guida EMA)
% -------------------------------------------------------
\begin{titlepage}
    \begin{center}
        \vspace*{1cm}
        {\large \textsc{Università degli Studi di Milano}} \\[0.4cm]
        {\normalsize Facoltà di Scienze Politiche Economiche e Sociali} \\[0.2cm]
        {\normalsize Corso di Laurea in Economia e Management}

        \vspace{3.5cm}

        {\LARGE \textbf{TITOLO DELLA TESI}}

        \vspace{4cm}

        \begin{flushleft}
            {\normalsize \textbf{Relatore:} Prof.\ Nome Cognome} \\[0.3cm]
            {\normalsize \textbf{Correlatore:} Prof.\ Nome Cognome}
                % Rimuovere la riga sopra se non presente
        \end{flushleft}

        \vfill

        \begin{flushleft}
            {\normalsize Tesi di Laurea di:} \\[0.2cm]
            {\normalsize \textbf{Nome Cognome Studente}} \\
            {\normalsize Matr.\ XXXXXX}
        \end{flushleft}

        \vspace{1cm}

        {\normalsize Anno Accademico 20XX\,/\,20XX}
    \end{center}
\end{titlepage}

% -------------------------------------------------------
% PAGINA BIANCA dopo il frontespizio (opzionale)
% -------------------------------------------------------
\thispagestyle{empty}
\cleardoublepage

% -------------------------------------------------------
% INDICE
% (inizia su pagina dispari, numerazione in romano)
% -------------------------------------------------------
\pagenumbering{roman}
\tableofcontents
\cleardoublepage

% -------------------------------------------------------
% TESTO PRINCIPALE
% (numerazione araba dalla prima pagina di testo)
% -------------------------------------------------------
\pagenumbering{arabic}

% Introduzione: non numerata, ma presente nell'indice
\chapter*{Introduzione}
\addcontentsline{toc}{chapter}{Introduzione}
\markboth{INTRODUZIONE}{INTRODUZIONE}
\begin{itemize}
    \item Il conflitto tra demografia e sostenibilità finanziaria.
    \item Obiettivi della tesi: analizzare le criticità storiche per delineare strategie future.
\end{itemize}
   % usare \chapter*{Introduzione} + \addcontentsline

% Capitoli numerati
%% ───────────────────────────────────────────────────────────────────
\chapter{Gli Obiettivi di un Sistema Pensionistico: Quadro Teorico}
%% ───────────────────────────────────────────────────────────────────

%% ─── PROMPT LLM ─────────────────────────────────────────────────
%% Scrivi il Capitolo 1 di una tesi magistrale in economia pubblica.
%% Obiettivo: fornire il framework teorico che guiderà tutta la tesi.
%% Tono: rigore scientifico + prosa scorrevole. Citare Bosi e Artoni
%% come fonti teoriche primarie, Rosen-Gayer-Rapallini (McGraw-Hill 2023)
%% come manuale di riferimento del corso.
%%
%% IMPORTANTE: questo capitolo deve stabilire il "trilemma previdenziale"
%% (sostenibilità / adeguatezza / equità) come framework analitico
%% che verrà usato in tutta la tesi per valutare riforme e strategie.
%% ────────────────────────────────────────────────────────────────

\section{Le tre funzioni del sistema pensionistico pubblico}

    %% PROMPT: Spiega le tre funzioni (previdenziale, assicurativa,
    %% assistenziale) con definizioni precise da manuale.
    %% Per ciascuna: definizione tecnica + esempio concreto + quale
    %% tipologia di pensione vi corrisponde (vecchiaia, invalidità,
    %% sociale, superstiti). Concludi mostrando che il settore privato
    %% può svolgere queste funzioni solo parzialmente e imperfettamente
    %% (risparmio individuale, famiglia, volontariato) -- giustificando
    %% così l'intervento pubblico.
    \subsection{Funzione previdenziale: garantire un tenore di vita
                simile a quello dell'ultima fase lavorativa}
    \subsection{Funzione assicurativa: protezione dal rischio
                di riduzione o fluttuazione del reddito nel tempo}
    \subsection{Funzione assistenziale: reddito minimo
                per un'esistenza dignitosa}
    \subsection{Il privato può sostituire il pubblico?
                Risparmio individuale, famiglia e volontariato
                come alternative imperfette}

\section{Le ragioni dell'intervento pubblico obbligatorio}

    %% PROMPT: Spiega le quattro giustificazioni dell'intervento
    %% pubblico OBBLIGATORIO e UNIVERSALE (non solo facoltativo).
    %% 1. Fallimenti del mercato assicurativo privato:
    %%    - asimmetrie informative sul rischio degli investimenti
    %%    - assenza di copertura dal rischio inflazione
    %%    - assenza di economie di scala (costi di marketing/gestione)
    %% 2. Miopia individuale (bene meritorio / paternalismo):
    %%    gli individui sottostimano il rischio di povertà in vecchiaia
    %%    e risparmiano meno del socialmente ottimale.
    %% 3. Azzardo morale: senza obbligo, un individuo potrebbe
    %%    non risparmiare sapendo che lo Stato interverrà comunque.
    %% 4. Selezione avversa: senza universalità, solo i lavoratori
    %%    meno produttivi rimarrebbero nel sistema, riducendo il monte
    %%    contributivo e portando alla crisi del sistema.
    %% Fonti: Bosi, Artoni, Rosen-Gayer-Rapallini cap. 12.
    \subsection{Fallimenti del mercato assicurativo privato:
                selezione avversa, asimmetrie informative
                e assenza di economie di scala}
    \subsection{Miopia individuale e bene meritorio:
                il paternalismo come giustificazione
                dell'obbligatorietà}
    \subsection{Azzardo morale e universalità:
                perché senza contribuzione obbligatoria
                il sistema si sgretola}
    \subsection{Il fallimento informativo come giustificazione
                aggiuntiva: perché la maggioranza dei lavoratori
                non conosce il proprio tasso di sostituzione atteso}

\section{I modelli di finanziamento: PAYG vs Capitalizzazione}

    %% PROMPT: Presenta i tre modelli con rigore tecnico.
    %% Per ciascun modello: meccanismo di funzionamento,
    %% vantaggi, svantaggi, rischi specifici.
    %%
    %% MODELLO 1 -- Capitalizzazione (FF: Fully Funded):
    %%   ottica assicurativa individuale, rendimento = r del mercato.
    %%   Pensione: P_FF = c * S_t * (1+r)
    %%   Rischio: volatilità dei mercati finanziari.
    %%
    %% MODELLO 2 -- Ripartizione (PAYG: Pay-As-You-Go):
    %%   patto intergenerazionale, rendimento implicito = (1+m)(1+n).
    %%   Pensione: P_PAYG = c * S_t * (1+m)(1+n)
    %%   Rischio: transizione demografica (riduzione di n).
    %%   I due sistemi danno la stessa pensione se (1+r) ≈ (1+m+n)
    %%   [Teorema di Aaron].
    %%
    %% MODELLO 3 -- NDC (Notional Defined Contribution):
    %%   ibrido: PAYG nella sostanza, DC nella forma.
    %%   Il conto è "virtuale" (nessuna riserva reale),
    %%   l'interesse implicito g = media mobile 5 anni del PIL nominale.
    %%   Spiega perché è una "terza via" e i suoi vantaggi
    %%   rispetto ai due modelli puri.
    %%
    %% Aggiungi: gli effetti economici della previdenza PAYG sul
    %% comportamento individuale (effetto sostituzione della ricchezza,
    %% effetto anticipo pensionamento, effetto eredità).
    \subsection{Il sistema a capitalizzazione (FF):
                ottica assicurativa individuale e rischio finanziario}
    \subsection{Il sistema a ripartizione (PAYG):
                il patto intergenerazionale e il teorema di Aaron}
    \subsection{Il Notional Defined Contribution (NDC):
                la terza via ibrida -- PAYG nella sostanza,
                DC nella forma}
    \subsection{Gli effetti economici della previdenza PAYG:
                spiazzamento del risparmio privato, effetto anticipo
                pensionamento ed effetto eredità}

\section{Il trilemma previdenziale: il framework della tesi}

    %% PROMPT: Presenta il "trilemma sostenibilità–adeguatezza–equità"
    %% come il prisma analitico che guiderà tutta la tesi.
    %% Definisci con precisione i tre vertici:
    %%
    %% 1. SOSTENIBILITÀ FINANZIARIA:
    %%    equazione macroeconomica estesa:
    %%    S_p/PIL = (N_p/N_a) × (N_a/POP) × (POP/N_l) × [(S_p/N_p)/(PIL/N_l)]
    %%    Spiega i quattro fattori: istituzionale, demografico,
    %%    occupazionale, redistributivo (tasso di sostituzione medio).
    %%
    %% 2. ADEGUATEZZA:
    %%    tasso di sostituzione individuale (pensione / ultimo salario),
    %%    soglia di povertà relativa in età anziana,
    %%    concetto di "tenore di vita dignitoso" (funzione previdenziale).
    %%
    %% 3. EQUITÀ INTERGENERAZIONALE:
    %%    tasso di rendimento implicito (IRR) per generazione,
    %%    confronto Baby Boomer vs Millennials/Gen Z,
    %%    concetto di neutralità attuariale.
    %%
    %% Concludi spiegando la TESI CENTRALE: le riforme Dini e Fornero
    %% hanno risolto il vertice della sostenibilità ma hanno creato
    %% un problema crescente sui vertici adeguatezza ed equità.
    \subsection{Sostenibilità finanziaria: la decomposizione
                dell'equazione $S_p/PIL$ nei quattro fattori}
    \subsection{Adeguatezza delle prestazioni: tasso di sostituzione,
                pensione minima e povertà relativa in età anziana}
    \subsection{Equità intergenerazionale: il tasso di rendimento
                implicito (IRR) e la neutralità attuariale}
    \subsection{Il trade-off al centro della tesi: sostenibilità
                comprata al prezzo dell'adeguatezza futura}
     % usare \chapter{Titolo}
%% ───────────────────────────────────────────────────────────────────
\chapter{Il Sistema Italiano: Dalle Origini al NDC}
%% [SINTESI STORICA -- ~12-15 pagine come da indicazione del relatore]
%% ───────────────────────────────────────────────────────────────────

%% ─── PROMPT LLM ─────────────────────────────────────────────────
%% Scrivi il Capitolo 2 di una tesi magistrale in economia pubblica.
%% ATTENZIONE: il relatore ha indicato che questa parte deve essere
%% SINTETICA (~12-15 pagine). Non è un capitolo storico, è un capitolo
%% di CONTESTO. Ogni sottosezione deve essere funzionale a spiegare
%% il PERCHÉ delle criticità attuali, non a descrivere eventi storici.
%% Tono: denso, essenziale, zero verbosità.
%% Fonti principali: MEF/RGS, Fornero (2018), Lavoce.info,
%%                  Rosen-Gayer-Rapallini cap. 12.
%% ────────────────────────────────────────────────────────────────

\section{Dalle origini alla crisi del sistema retributivo}

    %% PROMPT: Spiega l'evoluzione storica in modo funzionale al
    %% problema attuale. Segui questo filo:
    %% - Perché si è scelto il PAYG nel dopoguerra? (necessità di
    %%   pagare subito le pensioni senza accumulare riserve)
    %% - Il metodo retributivo (DB, 1970-1995): cos'è, perché era
    %%   generoso (pensione = 2% × anni × retribuzione finale),
    %%   perché ha creato insostenibilità latente (le promesse erano
    %%   finanziate dal debito implicito, non dai contributi).
    %% - La crisi: mostra la correlazione storica riforme-debito/PIL
    %%   (grafico Favero, Lavoce.info 2018): ogni sistema generoso
    %%   si è associato ad aumento del debito pubblico.
    %% - Spiega la pressione demografica usando la formula:
    %%   P_{t+1} = c*S_t*(1+m)*(1+n) → quando n ↓, le pensioni
    %%   diventano insostenibili senza ridurre P o aumentare c.
    \subsection{La scelta del PAYG nel dopoguerra:
                ragioni politiche ed economiche}
    \subsection{Il metodo retributivo (DB, 1970--1995):
                meccanismo, generosità e insostenibilità latente}
    \subsection{La pressione demografica e fiscale:
                l'invecchiamento nell'equazione $P = cS(1+m)(1+n)$}
    \subsection{La correlazione storica riforme--debito pubblico:
                1960--2018 come argomento per l'urgenza di riformare}

\section{La stagione delle riforme strutturali (1992--2011)}

    %% PROMPT: Tratta le tre riforme principali in modo comparato.
    %% Per ciascuna: contesto, misura principale, logica economica,
    %% impatto su sostenibilità e adeguatezza.
    %%
    %% RIFORMA AMATO (1992): primo intervento di emergenza.
    %%   - Modifica al calcolo retributivo (da 5 a 10 anni per i privati)
    %%   - Innalzamento età pensionabile (da 60 a 65 anni uomini)
    %%   - Contesto: crisi valutaria, pressione UE
    %%
    %% RIFORMA DINI (1995): la svolta strutturale.
    %%   - Introduce il NDC per i nuovi entrati dal 1/1/1996
    %%   - Il regime misto pro-rata per chi aveva < 18 anni di anzianità
    %%   - La logica: legare la pensione ai contributi versati
    %%     e alla crescita del PIL, non alle ultime retribuzioni
    %%   - Problema: effetti a regime solo dopo 30-40 anni
    %%
    %% RIFORMA FORNERO (2011): completamento e urgenza.
    %%   - NDC per tutti per gli anni residui (anche il regime misto)
    %%   - Aggancio automatico dell'età pensionabile alla speranza di vita
    %%   - Contesto: crisi spread 2011, pressione BCE/FMI
    %%   - Paradosso: la riforma "salva i conti" ma crea il problema
    %%     delle pensioni future inadeguate per i giovani
    %%
    %% Costruisci una tabella comparativa delle tre riforme:
    %% | Riforma | Anno | Misura chiave | Impatto sostenibilità |
    %% |         |      |               | Impatto adeguatezza   |
    \subsection{La Riforma Amato (1992):
                primo intervento di emergenza sul retributivo}
    \subsection{La Riforma Dini (1995):
                introduzione del NDC e logica del montante contributivo}
    \subsection{La Riforma Fornero (2011):
                NDC per tutti e aggancio automatico alla speranza di vita}
    \subsection{Il regime misto ``a tre velocità'':
                retributivo puro, misto, contributivo puro --
                tre generazioni, tre destini pensionistici diversi}

\section{L'architettura attuale: meccanismi e anomalie}

    %% PROMPT: Descrivi il sistema come funziona OGGI (2026).
    %% Tre blocchi:
    %%
    %% BLOCCO 1 -- Meccanismo NDC in dettaglio:
    %%   - Aliquota contributiva: 33% dipendenti, 20% autonomi
    %%   - Interesse implicito g = media mobile quinquennale PIL nominale
    %%   - Coefficiente di trasformazione: valori 2024 (4,270% a 67 anni
    %%     fino a 6,655% a 71 anni), aggiornamento biennale
    %%   - Come si calcola concretamente la pensione (formula)
    %%
    %% BLOCCO 2 -- Le deroghe politiche:
    %%   - Quota 100 (2019-2021): 62 anni + 38 contributivi
    %%   - Quota 102 (2022), Quota 103 (2023-in corso)
    %%   - Perché queste deroghe violano la logica NDC: vanno in pensione
    %%     prima che il montante sia sufficiente, riducendo la pensione
    %%     propria E generando costo per il sistema
    %%   - Costo stimato di Quota 100: ~8 miliardi/anno (INPS)
    %%
    %% BLOCCO 3 -- I numeri della spesa:
    %%   - 15,4% PIL nel 2024, 37,2% di tutte le uscite correnti (UPB 2025)
    %%   - La "gobba" MEF nel 2044 (picco per pensionamento Baby Boomers)
    %%   - La discesa strutturale post-2044 per entrata a regime del NDC
    %%   - Gli scenari demografici ISTAT integrati nelle proiezioni
    \subsection{Come funziona il calcolo NDC oggi:
                aliquote, coefficienti 2024 e aggiornamento biennale}
    \subsection{Le deroghe politiche alla logica NDC:
                Quota 100, 102, 103 tra populismo e rigore attuariale}
    \subsection{La spesa pensionistica al 15,4\% del PIL (2024),
                il picco del 2044 e la discesa strutturale post-2044}
    \subsection{Il paradosso redistributivo:
                le pensioni riducono il Gini di $-10{,}7$ punti,
                più di IRPEF e trasferimenti in natura insieme
                (Banca d'Italia, 2025)}


\chapter{Oltre i Confini Nazionali: Modelli Comparati e Soluzioni Anticonvenzionali}
\label{chap:modelli}

\section{Diversificare il rischio: il modello Multipilastro (World Bank)}

\section{Analisi comparata internazionale}
\subsection{Il modello Svedese: trasparenza e "Busta Arancione"}
\subsection{Il modello a Capitalizzazione pura: il caso del Cile}

\section{Soluzioni Anticonvenzionali e Innovative}
\subsection{I Fondi Sovrani Pensionistici: il caso Norvegia è replicabile in Italia?}
\subsection{Fiscalità e tecnologia: l'ipotesi della "Robot Tax"}
\subsection{Politiche demografiche attive e immigrazione selettiva}

\section{L'Economia Politica della Riforma}
\subsection{Il "Transition Cost": il problema del doppio onere}
\subsection{Il paradosso dell'Elettore Mediano: perché è difficile riformare}

\chapter{Il Futuro nelle nostre mani: Strategie di "Self-Defense" Previdenziale}
\label{chap:futuro}

\section{La presa di coscienza per le nuove generazioni}
\subsection{Analisi del tasso di sostituzione atteso per Millennials e Gen Z}

\section{Il Secondo Pilastro in Italia}
\subsection{Funzionamento e vantaggi fiscali della previdenza complementare}
\subsection{TFR in azienda vs Fondo Pensione: un'analisi di convenienza}

\section{Gestione attiva del risparmio}
\subsection{L'interesse composto come alleato del lungo periodo}
\subsection{Il Capitale Umano come asset previdenziale (Lifelong Learning)}


% Conclusioni: non numerate, ma presenti nell'indice
\chapter*{Conclusioni}
\addcontentsline{toc}{chapter}{Conclusioni}
\markboth{CONCLUSIONI}{CONCLUSIONI}
    % usare \chapter*{Conclusioni} + \addcontentsline

% Bibliografia e Sitografia: non numerate
\chapter*{Bibliografia e Sitografia}
\addcontentsline{toc}{chapter}{Bibliografia e Sitografia}
\markboth{BIBLIOGRAFIA}{BIBLIOGRAFIA}
   % usare \chapter*{Bibliografia} + \addcontentsline

\end{document}
